%! Carrying Out a Test for the Difference of Two Population Means — Unit 7, Lesson 7.9 — Group Activity (Answer Key)
\documentclass[10pt]{article}
\usepackage{apstats-article}
\usepackage{apstats-key}   % pulls in -boxes; do NOT also load -boxes
\newcommand{\TallMath}[1]{$\displaystyle #1\rule[-1.4em]{0pt}{3.2em}$}

\begin{document}

\pageheader{Unit 7, Lesson 7.9}{Group Activity --- Answer Key}
\namepartnerperiod

\vspace{0.1in}

\textit{Complete all four SPDC steps for each scenario. Show all computations.}

\begin{scenariobox}[Tier R --- Remediate: Read, Interpret, and Conclude]{sky}
A food scientist tested whether a new recipe produces crackers with a different mean
sodium content than the original recipe. A two-sample $t$-test on independent random
samples gave: $t = 2.41$, $df = 58$, $p = 0.019$ (two-sided), $\alpha = 0.05$.

\begin{enumerate}[label=\alph*.]
  \item \textbf{State:} $H_0$: \ans{$\mu_1 - \mu_2 = 0$ (equal mean sodium)} \quad
        $H_a$: \ans{$\mu_1 - \mu_2 \ne 0$ (different means)}

  \item \textbf{Interpret the $p$-value.} Fill in the blanks.

        ``Assuming the two recipes produce crackers with \ans{equal} mean sodium content,
        there is a \ans{0.019} probability of observing a difference in sample means
        as \ans{extreme or more extreme} as the one found, just by chance.''

  \item \textbf{Decide.} At $\alpha = 0.05$: [\ans{reject} / fail to reject] $H_0$
        because $p = $ \ans{0.019} [\ans{is} / is not] less than $\alpha$.

  \item \textbf{Conclude.} Write a complete conclusion in context.
        \ansline{Because $p = 0.019 < \alpha = 0.05$, we reject $H_0$. We have convincing evidence that the two recipes produce crackers with different mean sodium content.}
        \ansline{}
\end{enumerate}
\end{scenariobox}

\begin{scenariobox}[Tier A --- Approaching: Compute and Conclude]{sky}
A sports physiologist wants to know whether a resistance training program increases mean
grip strength more than a stretching program. Participants were randomly assigned to one
of two groups. After 8 weeks: Group 1 (resistance), $n_1 = 24$, $\bar x_1 = 52.3$ kg,
$s_1 = 7.8$ kg; Group 2 (stretching), $n_2 = 26$, $\bar x_2 = 47.1$ kg, $s_2 = 8.4$ kg.
Dotplots are roughly symmetric with no outliers.

\begin{enumerate}[label=\alph*.]
  \item \textbf{State:} $H_0$: \ans{$\mu_1 - \mu_2 = 0$} \quad
        $H_a$: \ans{$\mu_1 - \mu_2 > 0$}
  \item \textbf{Plan:} Method: \ans{Two-sample $t$-test for a difference of means}. Conditions:
        \begin{itemize}
          \item Random: \ansline{Participants randomly assigned to groups --- randomized experiment.}
          \item Normal: \ansline{Both dotplots are roughly symmetric with no outliers; CLT applies ($n_1 = 24$, $n_2 = 26$).}
        \end{itemize}
  \item \textbf{Do:} $SE = \sqrt{\dfrac{s_1^2}{n_1} + \dfrac{s_2^2}{n_2}} = $ \ans{$\sqrt{7.8^2/24 + 8.4^2/26} \approx 2.23$}

        $t = \dfrac{(\bar x_1 - \bar x_2) - 0}{SE} = \dfrac{{\color{keyred}\mathbf{52.3 - 47.1}}}{{\color{keyred}\mathbf{2.23}}} \approx$ \ans{2.33}

        Technology: $p \approx$ \ans{0.012} \quad $df \approx$ \ans{47}

  \item \textbf{Conclude:} At $\alpha = 0.05$: \ansline{Because $p \approx 0.012 < \alpha = 0.05$, we reject $H_0$. We have convincing evidence that the resistance program increases mean grip strength more than the stretching program.}
  \ansline{}
\end{enumerate}
\end{scenariobox}

\begin{scenariobox}[Tier E --- Extension: One-Sided vs.\ Two-Sided]{sky}
Use the grip strength data from Tier A: $t \approx 2.33$, technology gives $df \approx 47$.

\begin{enumerate}[label=\alph*.]
  \item Suppose the physiologist had instead used $H_a: \mu_1 \ne \mu_2$ (two-sided).
        Use technology to find the two-sided $p$-value. \quad $p_{\text{two-sided}} \approx$ \ans{0.024}

  \item Compare the one-sided $p$-value (from Tier A) to the two-sided $p$-value.
        How do they relate to each other? \ansline{The two-sided $p$-value is exactly twice the one-sided $p$-value: $0.024 = 2 \times 0.012$.}

  \item The physiologist pre-specified $H_a: \mu_1 > \mu_2$ before collecting data.
        Is it appropriate to report the two-sided $p$-value instead because it yields
        a larger value? Explain. \ansline{No. The alternative hypothesis must be determined before looking at the data. Switching to a two-sided test after seeing the data is data snooping and inflates the Type I error rate.}
        \ansline{}

  \item Suppose the sample sizes were doubled to $n_1 = 48$, $n_2 = 52$, with the same
        sample means and standard deviations. Would the $SE$ increase, decrease, or stay
        the same? How would this affect $t$ and the $p$-value? \ansline{$SE$ would decrease (larger $n$ reduces variability). A smaller $SE$ gives a larger $|t|$, which produces a smaller $p$-value --- making it easier to detect a real difference.}
        \ansline{}
\end{enumerate}
\end{scenariobox}

\end{document}
